% ============================================================
%  Section 2 — Background
% ============================================================
\section{Background and Related Work}
\label{sec:background}

\subsection{LLM-Based Document Processing}

Early document processing pipelines relied on OCR plus rule templates — brittle,
format-specific, and expensive to maintain. The shift toward LLMs brought
zero-shot extraction quality competitive with trained classifiers~\cite{li2023structuredext},
but introduced new failure modes: hallucinated field values, unstable output format,
and insensitivity to company-specific constraints.

Recent benchmarks~\cite{invoicellm2024,ocrornot2025} show that \emph{native
image input} to multimodal LLMs often outperforms OCR-then-text pipelines,
especially when document layout carries semantic meaning. For Tunisian documents
specifically, the challenge is the French number format (\textit{1\,011,500}),
the mandatory Timbre Fiscal line, and the \textit{matricule fiscal} (7 digits plus
a letter) — none of which appear in the training corpora of general-purpose models.

Prompt design is the dominant lever. A study on Spanish electricity invoices found
that few-shot prompting outperformed zero-shot by 19 percentage points in
F1~\cite{electricityinvoice2025}, reaching 97.6\% with Gemini. This motivates our
use of SkillOpt to automate that prompt iteration.

\subsection{Retrieval-Augmented Generation}

RAG~\cite{lewis2020rag} conditions LLM generation on retrieved evidence. It can
reduce unsupported answers when retrieval succeeds, but it also introduces
retrieval errors and does not eliminate hallucinations. For account code assignment, the
knowledge base is small (20–50 rules per company) but the required precision is
high: one wrong code has direct financial consequences.

The standard pipeline — embed query, retrieve top-k, inject into prompt — has a
known weakness: dense retrieval struggles with exact keywords (account numbers,
supplier names). Modern production RAG stacks address this with:

\begin{itemize}[leftmargin=*, itemsep=2pt]
  \item \textbf{Hybrid search}: BM25 keyword search merged with dense vector
        search via Reciprocal Rank Fusion (RRF)~\cite{rrf2009}.
  \item \textbf{Reranking}: a cross-encoder (e.g., \texttt{bge-reranker-v2-m3})
        jointly scores query and each retrieved chunk for precision~\cite{bgem3}.
  \item \textbf{BGE-M3}~\cite{bgem3}: a single model that simultaneously supports
        dense, sparse, and multi-vector retrieval in 100+ languages, making it a
        natural upgrade for Tunisian (French/Arabic) documents.
\end{itemize}

\subsection{Multi-Agent LLM Systems}

Multi-agent frameworks such as AutoGen~\cite{autogen2023} and LangGraph decompose
complex workflows into specialized, independently-improvable agents. n8n adds a
visual orchestration layer with built-in HTTP retry and error branches — removing
the need to write orchestration code and making the workflow auditable at a glance.

\subsection{SkillOpt: Text-Space Prompt Optimization}

SkillOpt~\cite{skillopt2024} is a Microsoft Research system that treats the agent
skill file (system prompt) as a trainable parameter. The target LLM generates
rollouts on training examples; an optimizer LLM turns failed rollouts into bounded
add/delete/replace edits to the skill; a held-out validation gate accepts edits
only if they improve the score. Across six benchmarks and seven models, SkillOpt
achieves the best or tied-best result in all 52 evaluated cells, improving average
accuracy by +23.5 points over the no-skill baseline~\cite{skillopt2024}.

SkillOpt-Lite~\cite{skilloptlite2025} — a follow-up that formalizes skill
optimization as Zeroth-Order optimization — further reduces the number of rollouts
needed for convergence, achieving +25.4 points on GPT-5.4-nano with fewer steps.

\subsection{Tunisian Accounting Context}

The Tunisian \textit{Système Comptable des Entreprises} (SCE)~\cite{sce1997}
organizes expense accounts in Class~6 (codes 60x–65x). VAT rates are 7\%,
13\%, or 19\% depending on product category. Every commercial document must
include a Timbre Fiscal (fiscal stamp) line and the supplier's matricule fiscal.
DocDesk models both companies on the SCE structure; account numbers used here
are unverified against the official source.
